\documentclass[10pt, letterpaper]{article}

% --- PAKETI I PODEŠAVANJA ---
\usepackage[utf8]{inputenc}
\usepackage[serbian]{babel} 
\usepackage[margin=0.75in]{geometry} 
\usepackage{titlesec}
\usepackage{xcolor}
\usepackage{enumitem}
\usepackage{hyperref}

% --- DEFINICIJA BOJA ---
\definecolor{primary}{HTML}{111827}    
\definecolor{accent}{HTML}{E11D48}     
\definecolor{graytext}{HTML}{4B5563}   

% --- HIPERLINKOVI ---
\hypersetup{
    colorlinks=true,
    linkcolor=accent,
    urlcolor=accent,
    pdftitle={CV - Kosta Čolović}
}

% --- TIPOGRAFIJA I STRUKTURA ---
\urlstyle{same}
\pagestyle{empty} 

% Formatiranje naslova sekcija
\titleformat{\section}
  {\normalfont\Large\bfseries\color{primary}}
  {}{0em}
  {}
    [\color{accent}\titlerule \vspace{-1.4ex}]

\newcommand{\entry}[4]{
  \vspace{1.5ex}
  \noindent\textbf{#1} \hfill \textcolor{graytext}{\small #2} \\
  \noindent\textit{#3} \hfill \textcolor{graytext}{\small #4}
  \vspace{0.5ex}
}

% --- POČETAK DOKUMENTA ---
\begin{document}

% --- ZAGLAVLJE (Kontakt podaci) ---
\begin{center}
    {\Huge \bfseries \color{primary} Kosta Čolović} \\[1.5ex]
    \textcolor{graytext}{
        Beograd, Srbija \ $\bullet$ \ 
        +381 64 235 2533 \ $\bullet$ \ 
        \href{mailto:fkcolovic@gmail.com}{fkcolovic@gmail.com}
    } \\[0.8ex]
    \textcolor{graytext}{
        \href{https://kostaco.github.io}{Veb-sajt} \ $\bullet$ \ 
        \href{https://www.linkedin.com/in/kosta-\%C4\%8Dolovi\%C4\%87}{LinkedIn} \ $\bullet$ \ 
        \href{https://github.com/kostaco}{GitHub}
    }
\end{center}

\vspace{2ex}

% --- OBRAZOVANJE ---
\section*{Obrazovanje}
\entry{O.Š. Borislav Pekić, Beograd}{2013 -- 2021.}{Osnovno obrazovanje}{Prosek: 5.0/5.0}
\begin{itemize}[leftmargin=1.5em, topsep=0.2ex, itemsep=0.2ex]
    \item Đak generacije, Vukova diploma
    \item Brojne nagrade na takmičenjima iz matematike, fizike i biologije
\end{itemize}

\entry{Matematička gimnazija, Beograd}{2021 -- 2025.}{Srednje obrazovanje (učenici sa posebnim sposobnostima za matematiku)}{Prosek: 4.7/5.0}
\begin{itemize}[leftmargin=1.5em, topsep=0.2ex, itemsep=0.2ex]
    \item Maturski rad: \textit{Nadgledano mašinsko učenje} (nagrada "Najbolji maturski rad iz računarstva i informatike")
\end{itemize}

\entry{Matematički fakultet, Univerzitet u Beogradu}{2025 -- Danas}{Student osnovnih akademskih studija, smer Informatika}{}


\entry{Fizički fakultet, Univerzitet u Beogradu}{od oktobra 2026.}{Student osnovnih akademskih studija, smer Teorijska i eksperimentalna fizika}{}

\vspace{1ex}

% --- PROJEKTI ---
\section*{Projekti}

\entry{Ekonomija baze: Odabir optimalne osnove brojevnog sistema i primene}{2026.}{\href{https://kostaco.github.io/materijali/ekonomija-baze.pdf}{kostaco.github.io/materijali/ekonomija-baze.pdf}}{Istraživački rad}
\begin{itemize}[leftmargin=1.5em, topsep=0.2ex, itemsep=0.2ex]
    \item Analiza efikasnosti reprezentacije podataka kroz asimptotski dokaz teorijske optimalnosti baze $e$. Matematičko opravdanje praktične superiornosti ternarnog sistema (baza 3) uz prikaz primene u računarskoj arhitekturi i teoriji fraktala.
\end{itemize}

\entry{Contextual Timeline}{2026.}{\href{https://kostaco.github.io/contextual-timeline/}{kostaco.github.io/contextual-timeline}}{Web alat}
\begin{itemize}[leftmargin=1.5em, topsep=0.2ex, itemsep=0.2ex]
    \item Interaktivna platforma za vizuelizaciju i mapiranje istorijskih, naučnih i filozofskih događaja. Razvijeno korišćenjem modernih web tehnologija sa naglaskom na čist dizajn i intuitivan UI.
\end{itemize}


\entry{AI u astronomiji: Otkrivanje egzoplaneta i obrada svemirskih podataka}{2026.}{\href{https://kostaco.github.io/materijali/AIuAstronomiji.pdf}{kostaco.github.io/materijali/AIuAstronomiji.pdf}}{Istraživački rad}
\begin{itemize}[leftmargin=1.5em, topsep=0.2ex, itemsep=0.2ex]
    \item Pregled i poređenje klasičnih algoritama mašinskog učenja i dubokih neuronskih mreža za detekciju egzoplaneta nad velikim astronomskim skupovima podataka.
    \item Analiza primene AI modela u procesiranju svetlosnih krivih, spektroskopiji, visokokontrastnom snimanju i automatskom redukovanju šuma.
\end{itemize}


\entry{Nadgledano mašinsko učenje}{2024--2025.}{\href{https://kostaco.github.io/materijali/NadgledanoMasinskoUcenje.pdf}{kostaco.github.io/materijali/NadgledanoMasinskoUcenje.pdf}}{Istraživački rad}
\begin{itemize}[leftmargin=1.5em, topsep=0.2ex, itemsep=0.2ex]
    \item Nagrada "Najbolji maturski rad iz računarstva i informatike"
    \item Detaljan teorijski i praktični pregled algoritama nadgledanog učenja, sa fokusom na analizu regresije, stabala odlučivanja, SVM-a i neuronskih mreža.
\end{itemize}

\entry{AstroJump}{2023.}{\href{https://github.com/KostaCo/AstroJump}{github.com/KostaCo/AstroJump}}{Video-igra}
\begin{itemize}[leftmargin=1.5em, topsep=0.2ex, itemsep=0.2ex]
    \item Razvoj 2D \textit{endless-runner} igre u Unity okruženju uz implementaciju simulacije fizike kretanja i mehanike skokova.
    \item Proceduralno generisanje prepreka (asteroida) i arhitekturno upravljanje stanjima igre (\textit{Game Loop}) korišćenjem C\# skripti.
\end{itemize}

\vspace{2ex}

% --- ISTRAŽIVAČKA INTERESOVANJA ---
\section*{Istraživačka interesovanja}
\vspace{1ex}
\begin{itemize}[leftmargin=1.5em, topsep=0.8ex, itemsep=0.3ex]
    \item \textbf{Mašinsko učenje i primene}
    \item \textbf{Filozofija i veštačka inteligencija}
    \item \textbf{Informatika u fizici i astronomiji}
\end{itemize}

\vspace{2ex}

% --- VEŠTINE ---
\section*{Tehničke veštine}
\vspace{1ex}
\begin{itemize}[leftmargin=1.5em, topsep=2ex, itemsep=0.3ex]
    \item \textbf{Programski jezici:} C, C++, C\#, Python, ARM64 Assembly, x86-64 Assembly, JavaScript, MATLAB
    \item \textbf{Operativni sistemi i alati:} Linux, Git, \LaTeX, HTML, CSS, Unity
    \item \textbf{Strani jezici:} Srpski (maternji), Engleski (napredni nivo), Nemački (osnovni nivo), Latinski (osnovni nivo)
\end{itemize}

\vspace{2ex}

% --- ISKUSTVO ---
\section*{Radno iskustvo}
\entry{Servisiranje i optimizacija računarskih sistema}{2022--2025.}{Samostalni tehnički projekti i servis}{}
\begin{itemize}[leftmargin=1.5em, topsep=0.2ex, itemsep=0.2ex]
    \item Dijagnostika hardverskih kvarova, projektovanje i sklapanje namenskih radnih stanica, uz uspešno rešavanje softverskih konflikata na Linux i Windows platformama.
    \item Održavanje integriteta podataka pri migracijama sistema, konfiguracija sistemskog softvera i pružanje tehničke podrške za lokalne klijente tokom srednjoškolskog obrazovanja.
\end{itemize}

\end{document}